\subsection{A cyclic encoding of the partial paths}
\label{sec:projective-word}

Split the boundary into two labelled intervals of length \(n\), read both
intervals in increasing order, and let \(h\) be the number of chords of the
noncrossing matching that join the intervals. The \(h\) ends on one interval
are paired with those on the other in reverse order. We refer to the ends of
these prescribed chords as anchors. Locally paired ends are read with a stack.
An anchor can be consumed only when the local stack on its side is empty.
At each step the antipodal matching joins the two vertices just read. A
component closed before the last step can never be joined to a later vertex.

Write \(a,b\) for the numbers of unused anchors and \(\ell,r\) for the
numbers of local openings on the two sides. Order the free ports by taking the
left anchors in reverse order, the left local stack from bottom to top, the
right local stack from top to bottom, and the right anchors in forward order.
Draw a chord between the two ports of each partial path, and call two chords
crossing when their ends alternate. These chords record connectivity.

\paragraph{Lemma.}
The crossing graph of the partial paths is a clique together with isolated
vertices. On \(2m\) labelled ports all matchings of this form, including the
planar ones, admit a canonical binary encoding. Their number is
\[
 C_m+\sum_{k=2}^{m}\binom{2m}{m-k}.
\]

\paragraph{Proof.}
Suppose the clique has \(k\ge2\) chords. Every isolated chord cuts off a
unique cyclic interval containing no clique ends. Indeed, a clique chord
cannot have one end on each side of an isolated chord, and clique chords on
opposite sides could not cross. Orient the isolated chord along the interval
without clique ends. Mark its initial end with a one and its final end with a
zero. Mark every clique end with a one. Successive deletions of cyclically
adjacent \(10\) pairs recover the isolated chords from the inside out. The
remaining \(2k\) ends must be paired with their opposites in cyclic order.
Conversely, any binary word with \(m+k\) ones gives such a matching: the
cyclic cancellations leave precisely \(2k\) ones, and the recovered ordinary
arcs cut off no residual end. The cancellations are independent of the choice
of pair, since different currently cancellable pairs are disjoint. In the
planar case orient each chord from its smaller to its larger index to obtain
the usual rooted Dyck word.

Initially there are \(h\) mutually crossing chords, with the planar
convention when \(h=0,1\). Two openings insert \(10\) at the active cut.
One opening and one closing move the cut while preserving the word. Two
closings fuse adjacent ends. A pair \(10\) closes a component. For \(01\),
delete the bits. For \(00\), first change the opener of the first zero to a
zero. For \(11\), when the second one has an ordinary closing partner,
first change that partner to a one. These rules concatenate the two paths
and preserve the orientation on all other ordinary chords.
If the second one in \(11\) is a clique end, so is the first: an ordinary
opening there would enclose the second end. Their opposite partners are
consecutive among the remaining clique ends. Joining the partners makes an
ordinary chord and leaves \(k-2\) mutually crossing chords. If at most one
survives, recover the rooted planar encoding. This proves the invariant by
induction and the enumeration follows from the inverse construction.
\hfill\(\square\)

\paragraph{The state quotient.}
Set \(L=a+\ell,\ R=b+r\) and let \(w\) denote the word in the stated port
order. The triple \((L,R,w)\) suffices for every continuation. In zero-based
indices a left closing always takes port \(L-1\), and a right closing always
takes port \(L\). At an empty local stack the same index denotes the next
anchor. A closing on a side is available exactly when its total length is
positive. An opening adds one to this length. After the move both lengths
must fit within the remaining number of steps. At a successful terminal
both stacks and all anchors have been consumed. Consequently two prefixes
with different local and anchor counts, even from different initial
heights, have identical weighted continuations whenever their triples
agree. For each fixed starting height, every pair of legal step sequences
recovers a unique original diagram by stack matching and reverse gluing.

Exchanging the two sides reverses the cyclic order. It reverses the
orientation of ordinary chords and preserves the set of clique ends. It
interchanges the two mixed steps and preserves the terminal contribution.
Thus the entire accumulated weight of an orbit can be propagated through
all moves from any representative. Moves with coincident targets retain
their multiplicities. This calculation counts labelled diagrams.

Deleting the last two ports by a closure at time \(j\) contributes to the
count of order \(j\). Restricting the two intervals to the \(j\) vertices
already read preserves every antipodal pair in the relative labelling.
Hence a transfer with horizon \(n\) can accumulate all shorter terminal
counts of the same parity in its starting height interval. The size of
the crossing clique never increases, so the time evolution of all states
with a fixed clique size can also be completed before processing the
smaller size to which they contribute.

\subsection{An exponential bound for the exact transfer}
\label{sec:projective-work-bound}

\paragraph{Theorem.}
Let
\[
 S=\frac{25}{4},\qquad \lambda=\frac{25+\sqrt{561}}8,\qquad
 \alpha=\tfrac12\log_2\lambda,\qquad
 \gamma=2^{2\alpha/(2+\alpha)}=1.72772438\ldots .
\]
There is a deterministic algorithm for \(P_n\) with time and space
\(O(n^{O(1)}\gamma^n)\).

\paragraph{Proof.}
Count separately, for this bound, the states at time \(t\) with prescribed
\(h,\ell,r,a,b\). The number of such parameter choices is polynomial in
\(n\). Put
\[
 c_L=h-a,\quad c_R=h-b,\quad
 d_L=\ell+c_L,\quad d_R=r+c_R,\quad D=d_L+d_R,\quad
 s=\ell+r+a+b.
\]
With \(u=a+b-h\ge2\), at least \(u\) original paths have both ends still
untouched. Those \(u\) chords still cross pairwise. Their \(2u\) bits are
known ones in the two half-open intervals
\([h-b,a)\) and \([s-b,s+a-h)\). The residual bit length is \(s-2u=D\).
If \(u\le1\), then \(s=D+2u\le D+2\). There are therefore at most
\(4\,2^D\) states in the prescribed group. There are also at most \(2^s\),
where \(s\le2(n-t)\) because both stacks must be discharged.

Each side of a prefix is a sign walk starting at zero. Its minimum is
\(-c_L\) and its endpoint is \(\ell-c_L\) on the left, with the corresponding
formulas on the right. Reflect the signs of each walk up to its first
attainment of the minimum, inclusive. The transformed walks end at
\(d_L,d_R\). Recording the two reflection times makes this mapping
injective and partitions the steps into at most three intervals of
constant sign changes.

A retained prefix cannot have two simultaneous openings immediately
followed by two simultaneous closings: the latter would close the new
path to itself. Lando and Zvonkin used this forbidden two-letter factor
and the nonoverlapping words for irreducible closed systems to bound
the growth of the projective meandric numbers
\cite[Lemma~3.3 and Theorem~3.4]{lz}. Here we apply the restriction after
two changes of signs and weight it by final displacement to bound the
number of transfer states. On each transformed interval one consecutive
pair of letters \(X,Y\) is forbidden, where the four letters are
\((+1,+1),(+1,-1),(-1,+1),(-1,-1)\) and \(Y=-X\).
Weight a letter \((i,j)\) by \(2^{i+j}\). The total letter weight is \(S\),
and \(\mathrm{wt}(X)\mathrm{wt}(Y)=1\) for every possible orientation.
Let \(F_j\) be the total weight of length-\(j\) words avoiding this pair.
Appending any letter produces total weight \(S F_{j-1}\). The invalid
extensions are exactly a valid length-\((j-2)\) prefix followed by \(XY\),
since \(X\ne Y\). Thus
\[
 F_0=1,\qquad F_1=S,\qquad F_j=S F_{j-1}-F_{j-2}.
\]
Its largest characteristic root is \(\lambda\), giving
\(F_j\le C\lambda^j\) for an absolute \(C\). For fixed reflection times,
ignoring restrictions across the at most two interval boundaries gives
total weight at most \(C^3\lambda^t\). Each image contributing to the
prescribed group has weight \(2^D\). Including the at most \((t+1)^2\)
pairs of reflection times, the number of prefixes, and therefore of
states, is \(O((t+1)^2\lambda^t 2^{-D})\). Taking the geometric mean
with \(4\,2^D\) bounds each group by \(O((t+1)\lambda^{t/2})\).
Finally
\[
 \min\{\alpha t,2(n-t)\}\le\frac{2\alpha}{2+\alpha}\,n .
\]
A balanced search tree implements exact state aggregation in polynomial
overhead per transition. Every coefficient has \(O(n)\) bits: extend each
prefix deterministically to a full noncrossing matching, disregarding
connectedness, and recover the prefix from its completion and the cut.
The bound follows for both time and space. \hfill\(\square\)

\paragraph{A ballot check.}
Without excluding any local factor, reflection at the minimum shows that
the exact number of one-sided prefixes with the indicated \(c,\ell\) is
\[
 \binom{t}{(t+\ell+c)/2}-\binom{t}{(t+\ell+c)/2+1}.
\]
For \(x=D/(2t)\ge2/3\), the product of the two ballot bounds is at most
\(2^{2tH_2((1+x)/2)}\le 2^{2tH_2(5/6)}<2^{4t/3}\), where \(H_2\) is binary
entropy. For \(x\le2/3\) the bit-length bound gives \(4\,2^{4t/3}\).
Together with the remaining-length bound this yields the simpler estimate
\(O(n^{O(1)}2^{4n/5})\).

\subsection{Rooting at a minimum diameter cut}
\label{sec:minimum-cut}

Let \(q_t(A)\) be the number of chords of the noncrossing matching \(A\)
between the two open semicircles at the diameter cut moved \(t\) gaps from
its initial position. Cuts opposite each other coincide as unlabelled
diameters, so the indices form a cycle of length \(n\). At the initial
cut \(q_0=h\). During the simultaneous scan,
\[
 q_t=\ell+r+|a+b-h|.
\]
Indeed, the local arches crossing the moving cut contribute \(\ell+r\).
An original through chord contributes precisely when both of its ends have
been read, or neither has been read. Reverse ordering of the original
through endpoints makes these two counts \(\max(0,h-a-b)\) and
\(\max(0,a+b-h)\), respectively.

Keep only prefixes with \(q_t\ge h\). For a diagram with minimal cut at
zero let \(g_0(A)\) be the positive number of steps from the preceding
minimal cut to zero in the cyclic list of cuts. It equals \(n\) if zero is
the only minimum. The exact labelled count is
\[
 P_n=\sum_{\substack{A\ {\rm connected}\\q_0(A)=\min_jq_j(A)}}g_0(A).
\]
For a fixed labelled diagram, sum the corresponding preceding gaps over
all its minimal cuts indexed by \(0,\dots,n-1\). This sum is \(n\).
Next sum over all labelled diagrams. Rotating the labelling sends the
summand with a prescribed index \(j\) bijectively to the sum at index zero;
thus the double sum equals \(n\) times the right-hand side. This proof also
covers periodic diagrams, since the double sum uses the indices with their
full multiplicities.

At each state store \(c\), the number of prefixes, and \(d\), their sum of
distances since the most recent minimal cut. Start with \((1,0)\). On each
admissible nonterminal move propagate \((c,0)\) if the new cut is a minimum,
and \((c,d+c)\) otherwise. At the final closure add \(d+c\) to the answer.
Exchange of the two sides preserves both weights. This version retains the
two anchor counts, since the minimum constraint uses them. Fixed initial
heights index \emph{minimum-height} sectors here; these differ from the
through-height sectors at a prescribed cut in the unrestricted calculation.

\paragraph{Lemma.}
For \(n\ge3\) every noncrossing matching has a diameter cut with
\(q_t\le\lfloor n/3\rfloor\).

\paragraph{Proof.}
We apply the three-arc tree argument of Rabinovich and Raz
\cite[Propositions 5.1--5.2]{rabinovich} to the contour of the matching.
Take the plane tree dual to the \(n\) chords in the disk, using one vertex
per region. Its unit-speed contour \(f\) has period \(2n\), visiting the
regions at integer gaps. The number \(q_t\) is the tree distance from
\(f(t)\) to \(f(t+n)\): the unique tree path crosses precisely the chords
separating these two boundary gaps. Extend \(f\) linearly on each edge
traversal. The images of three equally spaced times have a median in the
metric tree. Each of the three intervening contour arcs meets this median,
since the image of that arc contains the tree path between its ends.
Choose one preimage from each arc.

Two of the three chosen preimages have circular distance at least
\(2n/3\). Otherwise, writing the three consecutive spacings around the
circle, each spacing would be either less than \(2n/3\) or greater than
\(4n/3\). They sum to \(2n\), so one must be greater than \(4n/3\).
Its complement would be an arc shorter than \(2n/3\) meeting all three
original thirds, which is impossible. The antipode of one of the two
preimages is therefore at circular distance at most \(n/3\) from the
other. Since \(f\) is 1-Lipschitz, some real cut has tree distance at most
\(n/3\). On each unit interval the two contour points traverse tree edges
at unit speed. Their distance is affine unless they meet while traversing
the same edge in opposite directions. In that case the endpoint distances
are one. As \(n/3\ge1\), an integer cut attains the claimed bound in both
cases. \hfill\(\square\)

For even \(n\), every connected diagram rooted at a zero cut has a unique
zero cut among the \(n\) possible diameters. Two distinct zero cuts would
separate a nonempty proper block of simultaneous steps whose vertices are
closed under both matchings. Equivalently, a second zero in the two
initially empty stacks would close an intermediate component. Thus the
minimum-zero sector contributes \(nM_{n/2}\), using the standard two-half
bijection for a prescribed zero cut. The standalone minimum-cut transfer
can calculate that sector from the empty word without external ordinary
meander values. The exact records identify the executor for each reported
sector.

At order 30 the minimum-cut transfer returned \(41\,646\,461\,012\,208\)
and visited \(23\,859\,763\) legal height-filtered transitions in the
original run, including attempts at heights subsequently ruled out by
the lemma. The unrestricted side-symmetrised transfer required
\(228\,059\,280\) transitions when summing the initial heights separately.
At orders 35, 38 and 39 the complete minimum-cut totals also agree with
the unrestricted transfer. These comparisons check the treatment of the
labelled root and terminal weights. They are separate from the independent
disk recursion through order 31.

A separate Catalan enumeration tested every one of the \(13\,402\,696\)
noncrossing matchings of orders 1 through 15. For each rotated cut it
measured the height directly by testing which chord endpoints belonged to
the semicircle. The bound \(\min_t q_t\le\lfloor n/3\rfloor\) held
for \(3\le n\le15\). The labelled gap sum held at every tested order,
including for connected matchings with nontrivial rotational stabilisers.
An additional finite audit of the exponential bound covered
all 83\,303 nonempty shape groups through order 20, comprising 1\,634\,252
distinct frontier states. The ballot estimates, known-bit masks, forbidden
two-step factor and reflected weighted-word bound were checked with
integer arithmetic.

\subsection{The distribution of unit cuts in odd order}
\label{sec:unit-cuts}

Let \(o_j\) count ordinary open meanders with \(j\) ordered crossings,
using the convention of \cite[\S1]{lz}, and put
\(O(z)=\sum_{j\ge1}o_jz^j\). In a pair of noncrossing arch diagrams for an
open meander the two unmatched ends are exposed. They are on opposite sides
for odd \(j\) and on a prescribed common side for even \(j\). Fixing the
other choice of side gives the same count by reflection.
For a connected projective diagram \(A\) of odd order \(n\), let \(m(A)\)
be the number of its \(n\) cuts of height one.

\paragraph{Theorem.}
For odd \(n\), the number with \(m(A)>0\) is
\[
 2n[z^n]\log(1+O(z))-n o_n .
\]
More generally, for \(r\ge2\),
\[
 \sum_A\binom{m(A)}r=\frac{2n}{r}[z^n]O(z)^r,
 \qquad \sum_A m(A)=n o_n .
\]
In particular all positive frequencies of the statistic \(m(A)\) are
determined by ordinary open meander numbers of orders at most \(n\).

\paragraph{Proof.}
At one prescribed unit cut, remove the unique chord crossing that cut.
In the two half intervals the resulting arches have one exposed unmatched
end on each side. Connectedness gives a single open meander, and the removed
chord is recovered uniquely. The count is \(o_n\).

Now prescribe \(r\ge2\) distinct cuts, with positive cyclic spacings
\(d_1,\ldots,d_r\), totalling \(n\). A chord cannot be the unique crossing
chord for two prescribed cuts. The two diameters divide the boundary into
four intervals. A chord crossing both joins a pair of opposite intervals;
all other chords would then be internal to one interval. The other pair of
opposite intervals is closed under both matchings, contradicting
connectedness.

Thus each prescribed cut has its own crossing chord. The marked diameters
partition the boundary into \(2r\) elementary intervals. As one moves from
interval to interval, the \(r\) semicircle-membership bits change singly in
cyclic order, then repeat in that order. Two different intervals have
membership vectors at Hamming distance one precisely when they are adjacent.
A crossing chord therefore joins the two intervals adjacent to one of the
two lifted gaps of its cut. All remaining chords are internal to one
elementary interval.

Pairing opposite elementary intervals gives one block per spacing. Each
block has two dangling ends after removal of the crossing chords, one for
each of its bounding cuts. Its internal arches together with antipodal
edges form a single open path covering all \(d_i\) crossings: any additional
cycle would remain a component of the original diagram. Conversely these
open paths can be assembled in cyclic order. Choose which lift of the first
cut contains its crossing chord. In a block of odd length the two exposed
ends must be on opposite sides, while at even length they must be on the
same side, since on each side the other vertices pair internally. Thus
each next lift is forced. After a circuit the side has switched because
\(n\) is odd, exactly as required by the antipodal identification. For a
specified choice of lifts there are \(o_{d_i}\) possibilities in block
\(i\). If both exposed ends lie on the same side, attach them in boundary
order. The connections across adjacent gaps are unique and noncrossing,
and join the open paths into a single cycle. The two initial lift choices
are distinct for \(r\ge2\), giving \(2\prod_i o_{d_i}\).

Summing over prescribed subsets gives the binomial moments: an ordered
composition and an initial cut have \(n\) placements, and each subset has
\(r\) choices of its initial cut. Inclusion--exclusion now gives the
logarithm. The full marked enumerator is equivalently
\[
 \sum_A\bigl(u^{m(A)}-1\bigr)
 =-n o_n(u-1)-2n[z^n]\log\!\left(1-(u-1)O(z)\right).
 \qquad\square
\]

For a division-free calculation write \(H_n=[z^n]\,zO'(z)/(1+O(z))\).
Then \(H_n=n o_n-\sum_{j=1}^{n-1}o_jH_{n-j}\), and the unit-cut
contribution is \(2H_n-n o_n\). Direct Catalan enumeration tested the
product formula separately for all 43\,682 nonempty marked subsets in odd
orders through 15. At order 43 the formula, using tabulated ordinary
open-meander numbers, and the direct minimum-cut transfer both give
\(125\,215\,915\,957\,522\,819\,244\) for this single sector.
All ordinary inputs required by the unit-cut formula through order 47
were also computed locally. A planar two-ended transfer gives
\(o_0,\ldots,o_{44}\); closing the two exposed ends gives
\(o_{45}=M_{23}\) and \(o_{47}=M_{24}\). The open channel of Osipov's
first-crossing code, at revision \texttt{be0a3f5}, gives
\(o_{46}=193\,909\,492\,888\,406\,631\,692\). The supplementary data
record the exact native count and its comparison with A005316.
